\documentclass[11pt,a4paper]{article}

\usepackage[utf8]{inputenc}
\usepackage[T1]{fontenc}
\usepackage[margin=1in]{geometry}
\usepackage{lmodern}
\usepackage{microtype}
\usepackage{booktabs}
\usepackage{tabularx}
\usepackage{array}
\usepackage{longtable}
\usepackage{multirow}
\usepackage{xcolor}
\usepackage{colortbl}
\usepackage{hyperref}
\usepackage{graphicx}
\usepackage{caption}
\usepackage{amsmath}
\usepackage{enumitem}
\usepackage{titlesec}
\usepackage{listings}
\usepackage{float}
\usepackage{fancyhdr}

\hypersetup{
  colorlinks=true,
  linkcolor=blue!50!black,
  citecolor=blue!50!black,
  urlcolor=blue!60!black,
  pdftitle={Thermal Nexus -- Work Completed},
  pdfauthor={Thermal Nexus Team}
}

\definecolor{codebg}{rgb}{0.96,0.96,0.96}
\definecolor{codekw}{rgb}{0.13,0.13,0.55}
\definecolor{codecomment}{rgb}{0.45,0.45,0.45}
\definecolor{codestring}{rgb}{0.55,0.0,0.0}
\lstdefinestyle{shellstyle}{
  backgroundcolor=\color{codebg},
  basicstyle=\ttfamily\small,
  keywordstyle=\color{codekw}\bfseries,
  commentstyle=\color{codecomment}\itshape,
  stringstyle=\color{codestring},
  frame=single,
  framerule=0.4pt,
  rulecolor=\color{black!30},
  breaklines=true,
  showstringspaces=false,
  numbers=left,
  numberstyle=\tiny\color{black!50},
  numbersep=6pt,
}

\titleformat{\section}{\large\bfseries\color{blue!50!black}}{\thesection}{0.6em}{}
\titleformat{\subsection}{\normalsize\bfseries\color{blue!40!black}}{\thesubsection}{0.6em}{}
\titleformat{\subsubsection}{\normalsize\bfseries\color{black!70}}{\thesubsubsection}{0.6em}{}

\title{\textbf{Thermal Nexus --- Work Completed}\\
\large IEEE HART HardwAIre Challenge 2026 \\ Software Phase Summary}
\author{Thermal Nexus Team \\ \texttt{Pshawon-kundu / THERMAL-NEXUS}}
\date{\today}

\pagestyle{fancy}
\fancyhf{}
\fancyhead[L]{\small Thermal Nexus --- Work Completed}
\fancyhead[R]{\small IEEE HART HardwAIre Challenge 2026}
\fancyfoot[C]{\thepage}
\renewcommand{\headrulewidth}{0.4pt}

\begin{document}
\maketitle
\thispagestyle{fancy}

\section*{Abstract}
This report summarises the work completed so far on the Thermal Nexus predictive
cold-chain monitoring system. The main software pipeline --- synthetic data
generation, three operating modes, end-to-end radio+reader simulation, SQLite
storage, an offline Streamlit dashboard, KPI and energy analysis, and an
embedded model-export workflow with golden-vector parity testing --- is fully
implemented and tested. Building on that foundation, the team integrated an
external building-environment T15 benchmark dataset as a separate validation
stream, audited and converted the four user-supplied external sources, expanded
the native training corpus to a reproducible 96-run / 12{,}584-row
custom dataset, and refined the embedded parity and policy layers. The current
full suite reports 65 passed, 6 expected skips, 0 failures on a clean clone.

\section{Core Thermal Nexus software --- complete}
The main software pipeline is fully implemented under the IEEE HART HardwAIre
Challenge 2026 Phase 1--6 software scope. The completed modules include:
\begin{itemize}[leftmargin=*]
  \item Synthetic temperature-data generation (\texttt{simulator/temperature/}).
  \item Three sensor-node operating modes: fixed, rule-based, and ML-based
        transmission (\texttt{simulator/run\_end\_to\_end.py}).
  \item Deterministic radio channel + reader-side data collection
        (\texttt{simulator/radio}, \texttt{simulator/reader}).
  \item SQLite database storage with v1 $\rightarrow$ v2 migration
        (\texttt{host/database/}).
  \item Offline Streamlit dashboard with eleven sidebar pages
        (\texttt{host/dashboard/}).
  \item KPI generation, mode comparison, and energy-consumption estimation
        (\texttt{analysis/kpi\_engine.py}, \texttt{kpi\_reports.py}).
  \item Alert and event logging linked to radio/reader events
        (\texttt{host/database/alerts}).
  \item Embedded model export to C99 (\texttt{embedded/deployment/}).
  \item Golden-vector generation (\texttt{embedded/deployment/generate\_golden\_vectors.py}).
  \item Embedded parity-testing workflow (\texttt{embedded/tests/test\_parity.py}).
  \item Deployment and hardware-integration documentation
        (\texttt{docs/SYSTEM\_ARCHITECTURE.md},
        \texttt{EMBEDDED\_DEPLOYMENT\_SPECIFICATION.md}).
\end{itemize}

\subsection{Current test suite status}
\begin{center}
\begin{tabular}{lcc}
\toprule
\textbf{Suite} & \textbf{Passed} & \textbf{Skipped} \\
\midrule
Full \texttt{pytest} (clean clone) & 65 & 6 \\
Embedded parity (C99 golden vectors) & 12 / 12 & 0 \\
Class-mismatch vs sklearn & 0 & -- \\
Uniform-probability stubs detected & 0 & -- \\
\bottomrule
\end{tabular}
\end{center}
The six expected skips are gated integration tests that require the
\texttt{NEW-DATA-*.T15.txt} raw files or \texttt{--runs-hardware} on the CLI,
both intentionally absent from a clean clone.

\subsection{Embedded parity numbers}
\begin{center}
\begin{tabular}{ll}
\toprule
\textbf{Quantity} & \textbf{Value} \\
\midrule
Compiler & \texttt{gcc -std=c99 -Wall -Wextra -Werror} \\
Golden vectors tested & 12 \\
Max probability error (Python vs C99) & $2.60 \times 10^{-7}$ \\
Class mismatches & 0 \\
Policy mismatches & 0 \\
Uniform-stub vectors & 0 \\
Estimated static RAM & 160 bytes \\
Estimated flash & 18{,}667 bytes \\
\bottomrule
\end{tabular}
\end{center}

\section{Custom THERMAL-NEXUS dataset --- created}
The user-supplied external sources in \texttt{Datasets/} did not match the
THERMAL-NEXUS feature/label contract. The team built a reproducible native
dataset as the primary training corpus, and added a separate converter for
external sources as auxiliary domain-shift evidence.

\subsection{Audit of supplied external sources}
\begin{center}
\begin{tabular}{lrrl}
\toprule
\textbf{Source} & \textbf{Files} & \textbf{Compatible} & \textbf{Use} \\
\midrule
ELM \texttt{optimizationComp\_data.csv} & 1 & 0 (binary target) & auxiliary \\
Nepal vaccine-carrier loggers & 526 & 303 & domain-shift only \\
Research PDFs (4) & 4 & 0 (no tabular data) & reference \\
\textbf{TOTAL} & \textbf{531} & \textbf{303} & -- \\
\bottomrule
\end{tabular}
\end{center}
The Nepal carriers report 25--30~$^{\circ}$C temperatures and therefore map
almost entirely to \texttt{EXCURSION\_RISK} under the 2--8~$^{\circ}$C
cold-chain limits. They are useful as deployment-distribution evidence but
not as training data.

\subsection{Native dataset composition}
The native corpus is generated by \texttt{simulator/temperature/generate.py}
from the 12 scenarios in \texttt{config/scenarios.yaml} with 8 deterministic
seeds per scenario.

\begin{center}
\begin{tabular}{ll}
\toprule
\textbf{Quantity} & \textbf{Value} \\
\midrule
Scenarios & 12 \\
Runs per scenario & 8 (deterministic) \\
Total runs & 96 \\
Sample interval & 60 s \\
Duration per run & 7{,}200 s \\
Rows per run & $\approx$121 \\
Model-ready rows & 12{,}584 \\
Feature-valid rows & 11{,}833 \\
\bottomrule
\end{tabular}
\end{center}

\subsection{Class balance}
\begin{center}
\begin{tabular}{lrr}
\toprule
\textbf{State} & \textbf{Rows} & \textbf{Share} \\
\midrule
\rowcolor{blue!8} STABLE & 6{,}048 & 48.1\% \\
\rowcolor{blue!4} EXCURSION\_RISK & 4{,}592 & 36.5\% \\
\rowcolor{blue!4} TRANSITION & 1{,}112 & 8.8\% \\
\rowcolor{gray!10} UNAVAILABLE (boundary) & 832 & 6.6\% \\
\bottomrule
\end{tabular}
\end{center}

\subsection{Leakage-safe splits}
\begin{center}
\begin{tabular}{lrrl}
\toprule
\textbf{Split} & \textbf{Runs} & \textbf{Rows} & \textbf{Path} \\
\midrule
Train & 78 & 8{,}173 & \texttt{Datasets/CustomDataset/splits/train.csv} \\
Validation & 13 & 1{,}354 & \texttt{Datasets/CustomDataset/splits/validation.csv} \\
Test & 13 & 1{,}367 & \texttt{Datasets/CustomDataset/splits/test.csv} \\
\bottomrule
\end{tabular}
\end{center}
Split policy: stratified group K-fold by \texttt{run\_id}; zero overlap
between splits.

\section{Why our dataset is better suited to this project}
\begin{center}
\begin{tabularx}{\textwidth}{l X X}
\toprule
\textbf{Aspect} & \textbf{CustomDataset} & \textbf{Supplied sources} \\
\midrule
Schema & Full THERMAL-NEXUS contract with timestamps, run IDs, limits,
validity flags, labels, and historical features. & ELM has a serial index and
binary target. Nepal files have inconsistent headers and no native state labels. \\
Labels & Three classes: STABLE, TRANSITION, EXCURSION\_RISK, generated by the
existing future-excursion policy at 5/10/15-minute horizons. & ELM has binary
\texttt{Critical\_Level}; Nepal has no THERMAL-NEXUS labels. \\
Coverage & 12 engineered scenarios including warming, cooling, door opening,
spikes, drift, missing samples, and sensor faults. & ELM represents one
uncontrolled sequence. Nepal is mainly warm tropical field data. \\
Split safety & Group-stratified by \texttt{run\_id}; no run overlap. & ELM is
one time sequence; supplied files were not prepared for this contract. \\
Reproducibility & Deterministic seed and one-command regeneration. & Field
logger timestamps and acquisition conditions are not reproducible. \\
Use & Primary controlled training corpus for this project's task. & Auxiliary
domain-shift and provenance evidence only. \\
\bottomrule
\end{tabularx}
\end{center}
The supplied data has not been discarded. It is preserved separately and
converted with provenance under \texttt{ml/data/external/thermal\_nexus/}.

\section{Training results --- refreshed native dataset}
All four candidate models were trained on the expanded native dataset.

\subsection{Validation comparison}
\begin{center}
\small
\begin{tabular}{lrrrrr}
\toprule
\textbf{Model} & \textbf{Macro F1} & \textbf{Excursion Recall} & \textbf{Transition Recall} & \textbf{Selection Score} & \textbf{Artifact (B)} \\
\midrule
\rowcolor{blue!8} Logistic regression (selected) & 0.7847 & 0.8464 & 0.6250 & 5.4020 & 10{,}874 \\
Decision tree & 0.8444 & 0.9393 & 0.5956 & 5.8604 & 13{,}989 \\
Random forest reference & 0.8637 & 0.9286 & 0.6103 & 5.9734 & 305{,}717 \\
Small MLP (best validation) & 0.8625 & 0.9339 & 0.5956 & 6.0175 & 63{,}098 \\
Fixed threshold baseline & 0.5855 & 0.8393 & 0.0000 & -- & 0 \\
Rule-based baseline & 0.5698 & 0.9071 & 0.1912 & -- & 0 \\
\bottomrule
\end{tabular}
\end{center}
The small MLP is the best validation reference, but the current C99 export
supports logistic regression. Therefore logistic regression is the selected
embedded-deployable model; the MLP and random forest remain reference models.

\subsection{Final test metrics: selected logistic regression}
\begin{center}
\begin{tabular}{ll}
\toprule
\textbf{Quantity} & \textbf{Value} \\
\midrule
Test rows & 1{,}367 \\
Excursion events & 12 \\
Detected & 12 \\
Missed & 0 \\
False-alert events & 3 \\
Unnecessary state changes & 66 \\
Macro F1 & 0.7912 \\
Weighted F1 & 0.8651 \\
Balanced accuracy & 0.8084 \\
Mean inference latency & 0.0007 ms \\
\bottomrule
\end{tabular}
\end{center}

\subsection{Per-class metrics}
\begin{center}
\begin{tabular}{lrrr}
\toprule
\textbf{Class} & \textbf{Precision} & \textbf{Recall} & \textbf{F1} \\
\midrule
STABLE & 0.8661 & 0.9061 & 0.8857 \\
TRANSITION & 0.5084 & 0.6691 & 0.5778 \\
EXCURSION\_RISK & 0.9794 & 0.8500 & 0.9101 \\
\bottomrule
\end{tabular}
\end{center}

\subsection{Embedded C99 parity}
\begin{center}
\begin{tabular}{ll}
\toprule
\textbf{Quantity} & \textbf{Value} \\
\midrule
Golden vectors & 12 \\
Max probability error & $2.60 \times 10^{-7}$ \\
Class mismatches & 0 \\
Uniform-stub vectors & 0 \\
Status & \texttt{pass} \\
\bottomrule
\end{tabular}
\end{center}

\section{External T15 benchmark --- integrated}
The external building-environment T15 dataset remains a separate benchmark.
It contains 4{,}137 rows across 45 runs, with 31 train runs, 6 validation runs,
and 8 test runs at a 15-minute sampling interval. It is explicitly labelled
\texttt{EXTERNAL DERIVED BENCHMARK} and is not claimed as cold-chain field
validation. Its temperature-only and multivariate model-ready datasets,
30/60/90-minute targets, reports, and limitations are retained under
\texttt{ml/data/external/t15/} and \texttt{evidence/external\_t15/}.

\section{Scientific limitations}
\begin{itemize}[leftmargin=*]
  \item The CustomDataset is synthetic software data, not measured hardware data.
  \item The supplied Nepal data is useful for distribution analysis but is not
        representative of the 2--8~$^{\circ}$C training envelope.
  \item Energy numbers remain \texttt{ESTIMATED\_SOFTWARE\_VALUE} until the
        fabricated demonstrator is measured with a current probe.
  \item The selected C99 model is logistic regression. The higher-scoring MLP
        and random forest are reference candidates pending embedded exporters.
  \item One physical demonstrator is planned; multi-node scale remains simulated.
\end{itemize}

\section{Historical-only features and prediction horizons}
The main feature pipeline uses only current and historical measurements:
current temperature, lagged values, temperature differences and slopes,
rolling statistics, distance from limits, validity indicators, and gap
statistics. Future columns such as \texttt{future\_temperature\_*},
\texttt{will\_excursion\_*}, and \texttt{time\_to\_excursion\_*} are label
outputs and are excluded from predictors. The main model uses 5, 10, and
15-minute horizons. The external T15 benchmark separately uses 30, 60, and
90-minute horizons.

\section{Reports and reproducible workflow}
The following artefacts are available:
\begin{itemize}[leftmargin=*]
  \item \texttt{evidence/models/model\_comparison.csv}
  \item \texttt{evidence/models/final\_test\_metrics.json}
  \item \texttt{evidence/embedded/parity\_report.json}
  \item \texttt{evidence/embedded/resource\_estimate.json}
  \item \texttt{evidence/ai/energy\_policy\_summary.json}
  \item \texttt{simulation/radio\_link/results/}
  \item \texttt{scripts/prepare\_external\_dataset.py}
  \item \texttt{scripts/verify\_phase2.py}
  \item \texttt{docs/RUN\_MANUALLY.md}
\end{itemize}
The complete manual sequence is documented in \texttt{docs/RUN\_MANUALLY.md}:
generate data, build features, create group-safe splits, train candidates,
evaluate the selected model, export C99, run parity, run simulations, start
the dashboard, and execute the Phase-2 verifier.

\section{Clean-clone testing and test isolation}
The current full suite reports 65 passed, 6 expected skips, and 0 failures.
The six skips are integration tests gated on unavailable T15 raw data or
hardware. Test fixtures use temporary directories and no longer depend on
locally generated ignored files. Ruff and Black checks pass for the new
verification tooling.

\section{Phase-2 verification status}
\begin{center}
\begin{tabular}{ll}
\toprule
\textbf{Gate} & \textbf{Status} \\
\midrule
Embedded C99 compilation & PASS \\
Golden-vector compilation & PASS \\
Parity error $\leq 10^{-4}$ and no stubs & PASS \\
Embedded resource estimate & PASS \\
Energy-policy evidence & PASS \\
Equivalent engineering simulation & PASS \\
Full test suite & PASS: 65 passed, 6 skipped \\
Project description word limit & PASS: 1{,}049 words \\
Video script and shot list & PASS \\
Final video file & Deferred to Week 4 \\
Final submission ZIP & Deferred to Week 5 \\
\bottomrule
\end{tabular}
\end{center}

\section{Current project status}
\begin{center}
\begin{tabular}{ll}
\toprule
\textbf{Area} & \textbf{Status} \\
\midrule
Core simulator and reader pipeline & Complete \\
SQLite storage and migrations & Complete \\
Offline dashboard & Complete \\
KPI and energy analysis & Complete \\
Embedded export and parity & Complete \\
Custom dataset generation & Complete \\
External source audit and conversion & Complete \\
External T15 benchmark workflow & Complete \\
Feature engineering and leakage checks & Complete \\
Model training and evaluation & Complete \\
Scientific limitations & Documented \\
Manual run documentation & Complete \\
Phase-2 verification script & Implemented \\
Hardware measurements & Pending demonstrator \\
Final video and submission ZIP & Pending Week 4--5 \\
\bottomrule
\end{tabular}
\end{center}

\section{Conclusion}
The custom \texttt{Datasets/CustomDataset/} is the primary training corpus
because it is the only controlled source satisfying the THERMAL-NEXUS feature,
label, and split contracts. The supplied ELM and Nepal sources remain preserved
as auxiliary distribution evidence with explicit provenance. Training on the
expanded native dataset produces an embedded-deployable logistic-regression
model with 0 missed excursion events on 12 held-out events, 0.7912 test macro-F1,
and C99 probability parity within $2.60 \times 10^{-7}$. The project is
software-complete for the current phase; the remaining mandatory work is
physical demonstrator measurement, final video recording, and final submission
packaging.

\end{document}
